% !TeX root = ../../tesis.tex

\section{Apimetro: El motor de Datos metropolitano (Ingesta y ETL)}
\label{sec:apimetro-intro}

El proyecto \texttt{Apimetro} tiene su génesis entre los años 2023 y 2024,
derivado de una necesidad cartográfica y de modelación espacial para el universo
narrativo de \textbf{"La Saga de la Catacumba"}, obra del autor de la presente
investigación. El primer volumen de esta serie, titulado
\textbf{La historia de Quince}, explora una distopía urbana denominada «Ciudad
Capital», la cual funciona como un homónimo especulativo de la Ciudad de México
donde la habitabilidad se rige por la interacción crítica de tres variables:
transporte, educación y tiempo. En dicha narrativa, el «tiempo» no es solo una
dimensión física, sino un recurso de equidad social condicionado por la
eficiencia de la infraestructura de movilidad. Para la construcción de este
entorno literario, se requirió el desarrollo de un plan urbanístico que
resolviera la fricción espacial de los protagonistas, quienes residían en polos
opuestos de la metrópoli: el oriente y el sur profundo. En la realidad
geográfica de la Ciudad de México, estas zonas corresponden a demarcaciones con
un rezago histórico en transporte masivo, tales como Milpa Alta, Xochimilco,
Tláhuac y Tlalpan. En estos sectores, la oferta de movilidad recae
primordialmente en la \textbf{Red de Transporte de Pasajeros (RTP)}, cuya baja
frecuencia y velocidades comerciales limitadas generaban, a nivel narrativo,
«tiempos muertos» que dificultaban la fluidez de la trama. Esta problemática
orilló a una convergencia multidisciplinaria que dio origen a cuatro proyectos
paralelos: \textit{Apimetro}, \textit{La Catacumba}, \textit{VFTModel} y
\textit{Diarios}.

La primera aproximación técnica para sistematizar la geografía metropolitana fue
el desarrollo de \textit{«TrainMap»}, una herramienta de visualización que,
aunque funcional para propósitos creativos, carecía del rigor matemático
necesario para una auditoría urbana formal. A partir del año 2025, con el
ingreso a la \textbf{Maestría en Urbanismo de la UNAM}, la investigación dio un
giro hacia la formalización científica, integrando estándares como el
\textit{General Transit Feed Specification} (GTFS) y sistemas de información
geográfica (SIG).

\texttt{Apimetro} recibió una actualización estructural significativa a mediados
de marzo de 2026, orientándose hacia un modelo de
\textbf{Software as a Service (SaaS)} especializado en el manejo de tipos de
datos \texttt{GeoJSON} para la representación cartográfica y el análisis
ruteable. El núcleo operativo del sistema se define mediante un proceso
\textbf{ETL}
\footnote{En esta investigación se adopta la denominación \textit{Extraction, Treatment and Liberation} para enfatizar el carácter transformativo y de acceso abierto del proceso; el acrónimo estándar de la industria es «\textit{Extract, Transform, Load}».}
, diseñado para transformar la rigidez de los archivos \gtfs{} en geometrías
flexibles. Si bien el estándar \gtfs{} posee una estructura relacional robusta,
su formato original dificulta el análisis directo en hojas de cálculo o capas
\texttt{SHP} convencionales, por lo que la conversión a \texttt{GeoJSON} se
determinó como la solución óptima para la visualización y el cálculo de
indicadores topológicos. La persistencia y estructuración de la información se
ejecutan dentro de una base de datos relacional \textbf{PostgreSQL 15},
potenciada por la extensión espacial \textbf{PostGIS 3.4}. Esta elección
tecnológica permite almacenar datos geográficos de manera eficiente y
proporciona operaciones de consulta avanzadas, como \texttt{ST\_AsGeoJSON} y
\texttt{ST\_DWithin}, fundamentales para la ingesta y el ruteo asíncrono. Para
la implementación de esta fase, se utilizó el lenguaje de programación
\textbf{Python}, aprovechando la robustez de librerías especializadas en el
manejo de grandes volúmenes de datos como \texttt{Pandas} y \texttt{PySpark}.
Gracias a esta arquitectura, la ingesta de datos se consume de manera
autosuficiente y óptima, resultando en una base de datos de apenas 38 MB en
estado de desarrollo que es fácilmente replicable en diversos entornos
operativos, incluyendo Linux, Windows, MacOS y plataformas web. Esta
portabilidad facilita la auditoría de la red mediante gestores de consultas como
\texttt{Postman} o su despliegue local a través de tecnologías de contenedores
como \textbf{Docker} y \textbf{Docker Compose}. En última instancia,
\texttt{Apimetro} garantiza la soberanía tecnológica de la investigación,
permitiendo que la replicación del esquema de datos se realice de forma
transparente mediante el archivo \texttt{seed.sql} disponible en su repositorio
oficial.

\begin{figure}[htbp]
  \centering
  \includegraphics[width=0.45\textwidth]{Figures/Cap3/APImetroIcon.png}
  \caption[Ícono Personal de Apimetro]{Ícono Creado para el proyecto Apimetro}
  \label{fig:apimetro-icon}
  \fuentefigura{Elaboración propia.}
\end{figure}
